\documentclass{article}
\usepackage{pathfinder-paper}

\title{A Fixed-Price Baseline and a Rationing Boundary for Costly Auction Advice}

\begin{document}
\maketitle

\begin{abstract}
We examine the connection between a survey of learning-augmented algorithms and a local proposal for
token-metered language-model traders in a double auction. On the proposal's fixed value schedule, a scripted
price rule reaches the gross gains-from-trade ceiling of \(7.5\) if its five eligible pairs all execute. It
uses no generated language-model tokens, so paid decisions cannot improve net welfare over this completed
baseline on that instance. A second construction shows that an arbitrarily small change in an advised price
can permit a gross loss of \(1.25\) when arrival order rations a low-cost seller. Both statements are
conditional deductions, not reported auction outcomes. They motivate an information-matched,
allocation-aware test of optional price advice on varying value schedules.
% ledger 2, 10, 15, 16, 19, 20, 22
\end{abstract}

\section{The supplied pair and the question}

Q, the survey by Zhao et al.~\cite{zhao2026laa}, describes learning-augmented algorithms in terms of an
observable predictor input, a typed prediction, a task-specific error and a performance measure. It also
proposes charging predictor invocations in a cost-minimisation ratio against an uncharged optimum. These are
Q's framework and accounting convention, not auction guarantees.
% ledger 2, 7, 20, 21

P is the supplied local manuscript \emph{The Cost of Price Discovery in an LLM Double
Auction}~\cite{pathfinder2026p}, reproduced verbatim as Supplementary File S1
(\href{S1-P.tex}{S1-P.tex}) in the review materials. It proposes adding a generated-token charge to quote decisions in the
fixed-value auction of Struski et al.~\cite{struski2026market}. P derives the gross ceiling below and
proposes experimental controls, but reports no combined tariff experiment. P's own label ``Q'' refers to
Struski et al., not to the Zhao et al. survey assigned as Q here. The recorded search found no arXiv or DOI
record for P; the bibliography identifies the supplied manuscript by its supplementary-file identifier.
% ledger 1, 2, 7, 12

The new connection here is a precise boundary for optional price advice on P's instance. P meters tokens
generated by traders deciding quotes; it does not specify a separate predictor, prediction error or advice
interface. Consequently Q's consistency, robustness and charged competitive-ratio statements cannot be
transferred to P. We use an additive net-welfare comparison, and distinguish what follows from the value
schedule from what would require an execution or strategic model.
% ledger 2, 3, 5, 7, 8, 10, 22

\section{Accounting and a completed fixed-price comparator}

P's eleven buyers and eleven sellers have values and costs on the grid from \(0.75\) to \(3.25\) in steps of
\(0.25\). The five positive gains from efficient trades are \(2.5,2,1.5,1,0.5\); a sixth trade contributes
zero. Thus the maximum gross gain is \(G^{\star}=7.5\). Let \(G\) be realised gross gains from trade and
\(C\geq0\) the generated-token cost counted by P. With a declared conversion \(\rho\geq0\) from token cost
to value units, write
\[
W_{\rho}=G-\rho C,\qquad
R_{\rho}=G^{\star}-W_{\rho}=7.5-G+\rho C.
\]
This identity is accounting, not a learning-augmented guarantee. We use \(\rho\) because Q uses \(\lambda\)
for an unrelated trust parameter.
% ledger 3, 5, 7, 8, 16

\begin{proposition}[Conditional fixed-instance comparator]
On P's fixed schedule, let each buyer with \(v>2\) bid \(2+\varepsilon\), each seller with \(c<2\) ask
\(2-\varepsilon\), and all other traders abstain, where \(0<\varepsilon<0.25\). If the auction executes all
five eligible pairs, the script attains \(G=7.5\) and \(C=0\) under P's generated-token accounting. On this
completed path, every policy with \(C>0\) has strictly lower \(W_{\rho}\) when \(\rho>0\).
% ledger 4, 10, 14, 16
\end{proposition}

\begin{proof}
The five eligible buyer values sum to \(13.75\) and the five eligible seller costs to \(6.25\). Because
every eligible trader trades exactly once, pairing and execution price cancel from gross welfare:
\(G=13.75-6.25=7.5\). The quoted bid and ask cross, and \(\varepsilon<0.25\) keeps them inside the marginal
traders' values. The script generates no language-model tokens. No allocation can exceed \(G^{\star}\), so
any policy with \(C>0\) and \(\rho>0\) has \(W_{\rho}=G-\rho C<7.5\).
% ledger 10, 14, 16
\end{proof}

The execution premise is essential. P retains a random scheduler and a finite call cap; neither P nor the
ledger establishes that all five pairs trade when the script is used. This is neither an equilibrium nor a
claim that a strategic language-model trader obeys the script. In particular, the proposition does not
establish dominance over realised P runs or over variable value schedules.
% ledger 4, 10, 12, 14, 16, 22

\section{Price error and rationing}

Consider a stylised response in which buyers with \(v>p\) and sellers with \(c<p\) are eligible at advised
price \(p\). It makes the allocation issue visible without asserting how P's agents respond.
% ledger 15, 19, 22, 24

\begin{proposition}[Possible loss at arbitrarily small price error]
On P's same schedule, at \(p^{\star}=2\) there are five strictly eligible traders on each side. For every
\(0<\varepsilon<0.25\), at \(\hat p=2+\varepsilon\) the seller with cost \(2\) becomes eligible, while the
eligible buyer set stays fixed. A path that trades that seller and leaves out the seller with cost \(0.75\),
while executing five pairs, has \(G=6.25\), a gross loss of \(1.25\), even as \(\lvert\hat
p-p^{\star}\rvert\) tends to zero.
% ledger 15, 19, 22
\end{proposition}

\begin{proof}
The five original eligible sellers cost \(6.25\) in total. Replacing the seller costing \(0.75\) by the
seller costing \(2\) raises total matched cost to \(7.5\). The same five buyers total \(13.75\), hence
\(G=13.75-7.5=6.25\). This path can arise under an added price-response rule if the cost-\(2\) seller quotes
and a buyer crosses before the cost-\(0.75\) seller trades. It is a possible response and execution path,
not an observed run or a guarantee about P's implementation.
% ledger 15, 19, 22
\end{proof}

There is also a limited stability condition: if a price change leaves both eligible sets unchanged, each
contains \(K\) traders, and all \(K\) pairs execute, the gross sum is independent of price and pairing. On
P's schedule the type at \(2\) has zero threshold margin, so the eligible set can change after an
arbitrarily small upward error. The example therefore rules out a welfare-continuity claim based only on
scalar price error for this response and rationing model; it does not rule out a guarantee with additional
allocation assumptions.
% ledger 15, 19, 24

\section{Related work and interpretation}

Fixed-price double-auction mechanisms for gains from trade were already studied by Colini-Baldeschi et
al.~\cite{colini2017fixed}. The general step of using a fixed price to obtain high gains from trade is
therefore published prior work. Their distributional approximation analysis does not state the \(7.5\)
calculation for P's particular schedule or verify completion under P's random scheduler. Our comparator is
an instance-specific negative control, not a new general fixed-price mechanism. Clock auctions with
unreliable advice also have published welfare guarantees~\cite{gkatzelis2024clock}; their auction, advice
object and assumptions differ from the optional price consultation considered here. The targeted searches
recorded in \texttt{search.md} found no publication of this exact fixed-schedule rationing construction, but
do not establish novelty by absence.
% ledger 4, 10, 15, 19, 22

Q's cost term counts predictor invocations in a minimisation ratio; P maximises gross gains minus
generated-token cost. A trader-facing tariff \(\tau\) changes incentives, whereas \(\rho\) values the
resource cost for the analyst. Thus a future tariff comparison should evaluate \(G(\tau)-\rho C(\tau)\) at a
fixed declared \(\rho\). P already proposes declaring a range of welfare conversions, so this is an explicit
analysis distinction, not a demonstrated flaw in P.
% ledger 6, 7, 8, 9

\section{Limitations and open questions}

Neither proposition is an empirical outcome. The exact crossing, repeated-order and stopping behaviour of
the released auction must be checked before either possible path is asserted for that implementation. Finite
random scheduling can prevent complete matching, and strategic quote responses and incentive compatibility
are unproved. P's token accounting omits any resource consumed by a no-query script outside generated
language-model tokens. The posted-price response in the second proposition is an added model, not P's
specified agent policy.
% ledger 10, 12, 15, 19, 22, 24

The open empirical question is whether optional, paid price advice can improve net gains from trade on
prespecified varying value schedules when the efficient price is unknown. A test would give the adviser and
a no-query estimator the same observable private and public signals, declare the advice output and error
measure, keep scheduler and book rules fixed, and record gross gains, generated tokens, participation and
executions. P's fixed schedule provides a negative control under completed matching. Such a comparison could
answer that empirical question for its schedule family; it would not yield Q's general learning-augmented
guarantee.
% ledger 11, 12, 20, 21, 22, 24

\bibliographystyle{plainurl}
\bibliography{references}

\end{document}
